\chapter*{Preface}
Character sums of polynomials over finite fields give rise to error-correcting codes of remarkable
parameters. For a multiplicative character $\psi$ of $\bbF_q$ and polynomials $f \in \bbF_q[X]$ of
small degree, the codeword $(\psi(f(\alpha)))_{\alpha \in \bbF_q}$ has length $q$ and, by the Weil bound,
relative distance approaching $1/2$. Two specific families exhibit this behavior: the quadratic
residue character codes obtained by applying $\eta(\beta) = \beta^{(q-1)/2}$ to squarefree monic
polynomials of degree $O(\sqrt{q})$, and the dual BCH codes obtained over $\bbF_{2^b}$ by
composing the absolute trace with polynomials supported on odd-degree monomials. For both
families, the distance analysis reduces entirely to the Weil bound; in the case of dual BCH codes,
to such an extent that, absent the Weil bound, the existence of binary codes with these parameters
is not known by any other means. Until recently the decoding problem for these codes was open
even for non-constant degree; the situation was resolved by \cite{Kopparty2026}, whose
algorithm combines the Berlekamp--Welch like framework with the Stepanov polynomial method and
introduces a class of high-degree polynomials, pseudopolynomials, as the right algebraic object
for capturing the structure of character evaluations. The Stepanov method is also the elementary
route to the Weil bound itself; the same machinery that proves the bound thus also decodes the
codes it gives rise to. This report develops both threads in a unified, self-contained way.

The polynomial method of
Stepanov and Bombieri works directly via auxiliary polynomials on two specific curves:
$Y^d = f(X)$ for multiplicative character sums and $Y^q - Y = f(X)$ for additive character
sums. The primary references are \cite[Chapter~I-II]{Schmidt1976} and \cite[Chapter~2,5,6]{Lidl1996}.
The report is organized as follows. Chapter~1 develops the theory of finite fields: field extensions,
the Frobenius automorphism, and trace and norm maps. Chapter~2 treats characters of finite
abelian groups and finite fields, Gaussian and Jacobi sums, and develops the analytic infrastructure
of $L$-functions: the Riemann zeta function and its function-field analogue over $\bbF_q[X]$, the
Dirichlet $L$-function $L(s, \chi)$, which is then used to prove the Davenport--Hasse lifting theorem relating Gaussian
and Jacobi sums over $\bbF_q$ to those over its extensions. Chapter~3 addresses the problem of
counting solutions to equations over finite fields and derives Weil-type bounds for additive,
multiplicative, and quadratic character sums, together with the corresponding estimates for
quadratic forms and diagonal equations. Chapter~4 gives the two principal elementary proofs of
the Weil bound: Stepanov's proof for multiplicative character sums via the curve $Y^d = f(X)$,
and Bombieri's proof for additive character sums via the Artin--Schreier curve $Y^q - Y = f(X)$,
together with an alternate argument for the latter using B\'{e}zout's theorem. Chapter~5 develops
the theory of pseudopolynomials and the polynomial-time decoding algorithms for the quadratic
character and dual BCH code families, following Kopparty. Appendices~A--D collect supporting
material: algebraic and analytic prerequisites, modulus bounds from power sum estimates,
continued fractions over polynomial rings, and Hasse derivatives.

The report assumes familiarity with algebra at the level of groups, rings, and fields, and with linear
algebra. No prior exposure to algebraic geometry or analytic number theory is required.

\medskip
\begin{flushright}
Soham Chatterjee \\
School of Technology and Computer Science\\
Tata Institute of Fundamental Research, Mumbai \\
May 2026
\end{flushright}